\chapter[Mischehen zwischen Christen und Ungläubigen]{Was ist von Mischehen zwischen Christen und Ungläubigen zu halten?}
\quellseite{1}

1 Mos. 6, 2. berichtet der Heilige Geist von solchen Mischehen. Er sagt: \gesperrt{„Da sahen die Kinder Gottes“} (die zur Kirche gehörten) \gesperrt{„nach den Töchtern der Menschen“} (die ohne Gottes Wort dahinlebten), \gesperrt{„wie sie schön waren, und nahmen zu Weibern, welche sie wollten.“} Das war in den ersten Zeiten der Welt. Da fragten die jungen Männer nicht mehr darnach, ob die, die sie zu Weibern nehmen wollten, gläubig und gottesfürchtig waren. Nein, sie sahen nur darauf, ob es schöne Mädchen waren, die ihren Augen gefielen. Sie nahmen die Töchter der ungläubigen und gottlosen Weltkinder zu Weibern und ließen es sich nicht wehren. — So vermischte sich damals die Kirche Gottes mit der ungläubigen Welt.

Was sagte Gott der Herr dazu? „Da sprach der Herr: Die Menschen wollen sich von meinem Geist nicht mehr strafen lassen, denn sie sind Fleisch“, V. 3. Der Herr erkannte, daß die Menschen sich nicht mehr von seinem Geist durch sein Wort regieren lassen wollten, daß sie nicht mehr nach seiner Gnade und der ewigen Seligkeit fragten, daß sie nicht mehr göttlich leben wollten, sondern daß sie nach ihres verderbten Herzens Gelüste thun wollten, daß sie also vom Glauben abgefallen waren.

Und was war die Folge? Es erwuchs auf Erden ein gottloses Geschlecht, dessen Bosheit groß war, und alles Dichten und Trachten ihres Herzens war nur böse immerdar. Zwar waren gewaltige und berühmte Leute darunter. Aber was half das? Das war nur eine irdische Gewalt und Berühmtheit, V. 4. 5.

Und was war weiter die Folge? Die Sündfluth. — Warum erzählt der Heilige Geist das? Um uns zu warnen. Auch jetzt fragen viele junge Leute, die in der Kirche aufgewachsen sind, nicht mehr darnach, ob die, die sie heirathen wollen, gläubig und gottesfürchtig sind. Nein, sie fragen nur darnach, was ihnen gefällt, und nehmen ungläubige und gottlose Weltkinder zur Ehe und lassen es sich nicht wehren. — So vermischt sich auch jetzt die Kirche Gottes mit der ungläubigen Welt.

Was Gott der Herr jetzt dazu sagt? Ganz dasselbe, was er damals dazu gesagt hat. Auch jetzt wollen sich viele junge Leute nicht mehr von seinem Geist durch sein Wort regieren lassen, auch jetzt fragen sie nicht mehr nach Gottes Gnade und der ewigen Seligkeit, auch jetzt wollen sie nicht mehr gottselig leben, sondern sie wollen nach ihres verderbten Herzens Gelüste thun, auch jetzt sind viele vom Glauben abgefallen. Denn wenn das nicht so wäre, würden sie dann mit Ungläubigen und Gottlosen in die Ehe, in die allereinzigste Gemeinschaft, treten? Gewiß nicht.

Und was wird auch jetzt die Folge sein? Es wird ein gottloses Geschlecht erwachsen, dessen Bosheit groß und dessen Dichten und Trachten nur irdisch, weltlich, fleischlich, sündlich, böse sein wird immerdar. Zwar werden auch jetzt gewaltige und berühmte Leute darunter sein, Leute, die es auf dieser Welt zu etwas Großem bringen. Aber was hilft das?

Es gibt ja Ausnahmen. Es kommt ja durch Gottes Gnade vor, daß ein ungläubiger Mann oder ein ungläubiges Weib in der Ehe zur Kirche geführt und ein Kind Gottes wird. Aber in der Regel ist es doch so, wie der Heilige Geist sagt.

Und was wird jetzt weiter die Folge sein? Gottes Gericht. Ja, ja, wir sollen wissen, „daß der Himmel vorzeiten auch war, dazu die Erde aus Wasser und im Wasser bestanden durch Gottes Wort; dennoch ward zu der Zeit die Welt durch dieselbigen mit der Sündfluth verderbet. Also auch der Himmel jetzund und die Erde werden durch sein Wort gesparet, daß sie zum Feuer behalten werden am Tage des Gerichts und Verdammniß der gottlosen Menschen“, 2 Petr. 3, 5–7.

Darum sollen wir uns warnen lassen. —

Du fragst, was christliche Eltern und Pastoren denn machen sollen mit den jungen Leuten, die mit Ungläubigen Mischehen eingehen wollen? — Sie sollen sie warnen mit dem Gottesworte, das du eben gehört hast.

\quellseite{2}

Wird das helfen? — Wenn das junge Volk einmal los geworden ist, dann wird das Warnen wenig helfen. Die alten gläubigen Väter und Prediger vor der Sündfluth haben sicherlich auch gewarnt. Aber was hat es geholfen? Und zwingen — zwingen kann man nicht. Die Ursache sitzt inwendig, wie Gott spricht: „Die Menschen wollen sich meinen Geist nicht mehr strafen lassen, denn sie sind Fleisch.“ Aber dennoch — wir sollen warnen.

Aber ich will etwas sagen, was noch besser ist als warnen. Man thue allen Fleiß, die Kinder von Klein auf aufzuziehen in der Zucht und Vermahnung zum Herrn.

Ihr Eltern, laßt Gottes Wort in euren Häusern wohnen. Sagt euren kleinen und größeren und großen Kindern Gottes Wort. Haltet regelmäßige, kurze und einfache Hausandachten. Betet mit euren Kindern. Geht ihnen mit gutem Beispiel voran. Sendet eure Kinder in eine wahrhaft christliche Gemeindeschule. Wenn sie confirmirt sind, dann wacht über sie in der Liebe Christi. Betet für sie ohne Unterlaß.

Ihr Lehrer, haltet seine Religionsstunden. Macht diese den Kindern lieb und werth. Malt ihnen Jesum vor Augen. Lehrt sie die himmlische Lehre so, daß sie merken, daß ihr sie selig machen wollt.

Und ihr Pastoren, laßt euren ganz geringen Mitknecht euch ein Wort sagen. Haltet den Confirmandenunterricht und die Christenlehre und die Predigt ganz klar und kindlich und einfach, immer mit dem brünstigen Gebetsverlangen, daß dadurch der Geist Gottes in die Herzen komme.

Das sind die Mittel, die Kinder und die Jugend aufzuziehen in der Zucht und Vermahnung zum Herrn.

Das wird nicht, das kann nicht vergeblich bleiben.

Dadurch wird eine Jugend großgezogen, welche sich von Gottes Geist durch Gottes Wort regieren läßt, die nach Gottes Gnade und der ewigen Seligkeit fragt, die gottselig leben, die wider die verderbte Lust kämpfen und Glauben halten will.

Diese Jugend wird sich dann auch nicht mit der ungläubigen Welt vermählen wollen, sondern wird sich warnen lassen.

\begin{unterschrift}
E. M. B.
\end{unterschrift}
